% 实数（综述）
% license CCBYSA3
% type Wiki

本文根据 CC-BY-SA 协议转载翻译自维基百科\href{https://en.wikipedia.org/wiki/Real_number}{相关文章}

在数学中，实数是可以用来度量连续的一维量（如长度、时间或温度）的数。在这里，“连续”意味着数值对之间可以有任意小的差异。[a] 每一个实数都可以通过一个无限小数展开来几乎唯一地表示。[b]

实数在微积分（以及数学的许多其他分支）中是基础性的，尤其体现在其在极限、连续性和导数等经典定义中的作用。[c]

实数集有时被称为 “the reals”（实数），传统上记作黑体 R，通常使用黑板粗体符号 ⁠
$\mathbf{R}$⁠。[1][2] “实数” 这一形容词在 17 世纪由勒内·笛卡尔使用，用以区别实数与虚数（如 −1 的平方根）。[3]

实数包括有理数，例如整数 −5 和分数 $4/3$。其余的实数称为无理数。一些无理数（以及所有有理数）是具有整系数多项式的根，例如平方根 $\sqrt{2} = 1.414...$；这些数被称为代数数。而也存在一些并非如此的实数，例如 $\pi = 3.1415...$；这些数被称为超越数。[3]
